% -----------------------------
% 執筆にあたっての留意点
% -----------------------------
%
% 論文の第1章にあたる文章です．イントロダクションは極めて重要です．イントロダクションの役割は，読者に論文を読むべきか否かを敏速・的確に判断するための材料を提示すること，論文を読むために必要となる前提知識を提供することにあります．読者はイントロダクションを読んで面白いと思えば論文を読み続けるし，面白くなければ論文を読むのをやめます．イントロダクションで読者を惹きつけられなければ，そこで試合終了です．
%
% イントロダクションでは，以下の要素を順番に書いていくと良いでしょう．
%
% #### 手法・システム提案系論文の場合
% 1. （次に述べる問題につながる）背景
% 2. 何が問題なのか？
% 3. この論文で解決したいことは何か（目的）？
% 4. 掲げた目的が達成されると何が嬉しいのか（既存手法では何故ダメなのか）？
% 5. 提案手法・システムの直感的なアイデア
% 6. 提案手法・システムの具体的な説明
% （より詳細な説明は別の章で記すので，項目5を実現するための要素を書く程度でよい）
% 7. 研究で得られた知見（どんな実験をして，その結果，どんな知見が得られたか）
%
% 以下は例です．
% > 1. ウェブには正しい情報もあれば間違った情報もある．相当数の人が誤ったウェブ情報を鵜呑みにしている．
% > 2. どのウェブ情報が正しいか間違っているかについて，一般ユーザが常に注意を払うことは難しい
% > 3. 本稿では，信憑性が怪しい情報に対するユーザの注意力を高める方法を提案する
% > 4. 既存手法として，情報の信憑性を自動的に分析する方法が提案されているが，すべての情報を自動分析することは困難．最終的には人間に判断が委ねられるので，人間の力を高める必要がある．
% > 5. 積極的に証拠を調べる気にさせるメッセージをウェブ検索中に提示する
% > 6. 認知心理学のプライミング効果を利用して，批判的な情報の見方を促すメッセージを自動生成し，ウェブ検索時に表示する
% > 7. ユーザ実験の結果，通常のウェブ検索エンジンを用いるよりも，公的機関や学術機関が発行している情報を多く閲覧するようになった．
%
%
% #### ユーザ調査系の場合
% 1. （次に述べる問題につながる）背景
% 2. 何が問題なのか？
% 3. この論文で明らかにしたいことは何か（リサーチクエスチョンRQは何か）？
% 4. 掲げたRQが明らかになると何が嬉しいのか（既存研究の限界は何か）？
% 5. RQを明らかにするための調査・実験のアイデア
% 6. 提案する調査方法・実験方法の具体的な説明
% （より詳細な説明は別の章で記すので，項目5を実現するための要素を書く程度でよい）
% 7. 研究で得られた知見（調査・実験をした結果，どんな知見が得られたか）
%
% 以下は例です：
% > 1. ウェブには正しい情報もあれば間違った情報もある．
% > 2. 相当数の人が誤ったウェブ情報を鵜呑みにしている．
% > 3. どんなユーザが誤った情報を取得しやすいのかが分かっていない
% > 4. 情報を批判的に調べる能力についての調査は主に高等教育機関の学生を対象に行われているが，それ以外の対象に対する調査はない．またそれら調査は汎用的な認知能力（批判的思考能力）の調査が目的であり，ウェブ情報をどのように閲覧しているかの調査事例はほとんどない．
% > 5. アンケート調査と実際の検索行動分析を組み合わせた分析を行う．
% > 6. アンケート調査に回答した被験者に対して，ウェブ検索タスクをさせて行動を分析．
% > 7. 誤った情報を取得しやすいユーザは，積極的に情報を調べる意欲が乏しく，ウェブ検索をしても検索結果の上位しか閲覧しない．
%
% 各要素を書くときは，1段落で1主張となるようにすると読みやすくなります．
%
% 上記アドバイスを補足するとすれば，論文全体を通じて説明に用いる共通の例を導入することです．取り扱う問題やその重要性を分かりやすく伝えるために例を使うのは有効です．提案事項の詳細を述べる章でも同じ例を用いることで，読者は論文をスムーズに読み進めることができます．
%
%
% -----------------------------
% 補足（参考：how-to-xx/how-to-write-a-paper/1-学術論文の書き方.md 3.3節）
% -----------------------------
%
% ■ 行き詰まったら「5つの問い」に答えてみる
% 上の7項目と対応しますが，より短く覚えられる型として次のものがあります．書き出せなくなったら，この5問に一つずつ答えるつもりで書いてみてください．
%
% 1. What is the problem?（問題は何か）
% 2. Why is it interesting and important?（なぜ興味深く重要か）
% 3. Why is it hard?（なぜ難しいか．素朴なアプローチはなぜ失敗するか）
% 4. Why hasn't it been solved before?（なぜ未解決だったか．既存アプローチに何が足りないか）
% 5. What are the key components of my approach and results?（アプローチと結果の鍵は何か．限界も明記する）
%
% ■ 章の最後に「本論文の貢献」を箇条書きで置く
% イントロダクションの最後に，「本論文の貢献は次の3点である」と貢献を箇条書きで明示してください．その際，次の2点を守ります．
%
% * 各貢献は「検証可能な主張」として書く．読者が「本当か？」と確かめられる粒度にします．
%   （悪い例：「〜を改善した」／良い例：「〜を用いることで検索回数が有意に増加することを示した」）
% * 各貢献に，それを述べている章・節の番号を添える（例：「(2) 〜する手法を提案する（3章）」）
%
% 貢献に章番号を添えておけば，イントロダクション自体が論文全体の見取り図を兼ねます．その結果，
% 「2章では背景を述べ，3章では手法を述べ…」という中身のない章構成の羅列を書く必要がなくなります．
%
% ■ 査読者は，イントロダクションを読み終えた時点で採否の初期判断をほぼ下しています
% 残りの部分は，その判断を裏付ける証拠を探しながら読まれます．冒頭で「そこで試合終了」と書いたのは，この意味です．
%
% ■ 共通の例（running example）は，この章の末尾で導入します
% 導入位置は，イントロダクション末尾の小節とするか，独立した章（2章または3章）を立てるのが一般的です．
% 一度決めた例は使い捨てにせず，提案手法の章でも実験の章でも同じ例を育てていってください．

% ----------------------
% はじめに/introduction
% ----------------------
\section{はじめに}
